\section{Introduction}
\label{sec:intro}

\authornote{\textbf{P1 -- clinical motivation (MIDL requires it).} Accelerated MRI
shortens scans (motion, paediatrics, throughput). Undersampling makes reconstruction
ill-posed, and learned reconstructions can hallucinate structure. Cite the false
vessels and clefts reported in the 2020 fastMRI challenge
\citep{muckley2021fastmri} and the hallucination framework of \citet{bhadra2021hallucination}.
State the stakes: any change to the DC step changes what the reconstruction is allowed
to invent.}

\authornote{\textbf{P2 -- training-free diffusion priors.} One unconditional prior serves
many masks, accelerations and coil arrays \citep{jalal2021robust,chung2022scoremri,aali2025ambient},
unlike supervised unrolled networks \citep{sriram2020varnet}. Name the sampler families:
gradient guidance (DPS, \citealp{chung2023dps}), pseudo-inverse guidance
(\pigdm, \citealp{song2023pigdm}), Krylov/CG DC (DDS, \citealp{chung2024dds}),
plug-and-play (DiffPIR, \citealp{zhu2023diffpir}) and decoupled annealing
(DAPS, \citealp{zhang2025daps}). Common denominator: an \emph{isotropic} Gaussian
approximation of $p(x_0\mid x_t)$, as shown by \citet{peng2024optimal}.}

\authornote{\textbf{P3 -- the observation (teaser, \figureref{fig:teaser}).} With a fixed,
tuned configuration, PSNR \emph{decreases} as $T$ grows from \xx{} to \xx{} for \xx{} of
\xx{} samplers. Practitioners therefore have to tune $T$ per protocol. Related
evidence: SURE-based early stopping against noise overfitting in diffusion MRI
\citep{ozturkler2023smrd}, and the finite-sample analysis of miscalibrated
intermediate posteriors \citep{burns2026finite}. Make clear that this is the first
systematic step-count study on raw multi-coil MRI, and keep the claim exactly that
narrow.}

\authornote{\textbf{P4 -- key insight.} For Cartesian MRI the forward operator is diagonal
in k-space (up to coil sensitivities). Image power is far from white across spatial
frequency, while thermal noise after prewhitening is white
\citep{roemer1990phased,kellman2005snr}. The isotropic approximation therefore
over-trusts the measurements at frequencies where the signal is weak relative to the
noise, and under-trusts them where it is strong. A k-space-diagonal posterior variance
is the natural basis here. This contrasts with the wavelet basis of
\citet{peng2024optimal} and the dense Jacobian covariances of
\citet{boys2024tmpd,rozet2024mmps}.}

\authornote{\textbf{P5 -- contributions (3 bullets, each tied to a result).}
(i) A frequency-resolved DC update (\method) that is training-free, scan-calibrated
from the ACS region, closed-form for single-coil and CG-solved for SENSE, and plugs
into five samplers (\sectionref{sec:method}).
(ii) A step-count study on raw multi-coil brain and knee data, with a band-resolved
mechanism analysis (\sectionref{sec:res-steps}).
(iii) Validation on real low-field noise (M4Raw, 0.3\,T) and inside fastMRI+ pathology
regions (\sectionref{sec:res-snr}, \sectionref{sec:res-roi}).}

\begin{figure}[t]
\floatconts{fig:teaser}
{\caption{\textbf{More steps hurt isotropic data consistency.} PSNR vs.\ number of
reverse steps $T$ on raw multi-coil fastMRI brain ($R{=}\xx$, \xx{} test volumes,
median and IQR over volumes). Solid: samplers with isotropic DC, each tuned once on the
calibration split. Dashed: the same samplers with the frequency-resolved update.}}
{\figspec{4.2cm}{Line plot, x: $T\in\{10,20,50,100,200\}$ (log scale); y: PSNR (dB).
Lines: \ikdc, \pigdm, DDS, DiffPIR (solid) and +\method{} (dashed, same colour).
Shaded IQR bands. Source: \texttt{outputs/midl27/steps/*.json}.}}
\end{figure}
